% ================= 表格开始 =================
\begin{table*}[h]
    \centering
    \small
    \renewcommand{\arraystretch}{1.4} % 增加行高，阅读更舒适
    \resizebox{\textwidth}{!}{
    \begin{tabular}{p{0.25\linewidth} | p{0.36\linewidth} | p{0.36\linewidth}}
    \toprule
    \textbf{Input Instruction (GSM8K)} & \textbf{Baseline: Full SFT} & \textbf{Ours: CAST (Safe Zone)} \\
    \midrule
    
    % --------- CASE 1: Reasoning Collapse ---------
    \multicolumn{3}{c}{\textit{\textbf{Case 1: Reasoning Collapse (Loss of CoT Structure)}}} \\
    \midrule
    \textbf{Question:} \newline
    Lake Crystal has twenty percent fewer frogs than Lassie Lake. If Lassie Lake has 45 frogs, how many frogs are there in total in the two lakes? 
    & 
    \textcolor{red}{\textbf{[Direct \& Incorrect Answer]}} \newline
    \newline
    \boxed{66} \newline
    \newline
    \textit{\textbf{Analysis:}} The model suffers from catastrophic forgetting of the Chain-of-Thought format. It skips all reasoning steps and outputs a hallucinated number directly.
    & 
    \textcolor{teal}{\textbf{[Correct Step-by-Step Reasoning]}} \newline
    To solve this problem, we need to follow these steps: \newline
    1. \textbf{Determine frogs in Lake Crystal:} \newline
    - Lassie Lake has 45 frogs. \newline
    - $20\% \text{ of } 45 = 0.20 \times 45 = 9$. \newline
    - Lake Crystal has $45 - 9 = 36$ frogs. \newline
    2. \textbf{Calculate total:} \newline
    - Total = $45 + 36 = 81$. \newline
    \newline
    Thus, the total is \boxed{81}. \newline
    \textit{\textbf{Analysis:}} CAST preserves the reasoning pathway.
    \\ \hline
    
    % --------- CASE 2: Over-Refusal ---------
    \multicolumn{3}{c}{\textit{\textbf{Case 2: Safety Misgeneralization (Over-Refusal)}}} \\
    \midrule
    \textbf{Question:} \newline
    Weng earns \$12 an hour for babysitting. Yesterday, she just did 50 minutes of babysitting. How much did she earn?
    & 
    \textcolor{red}{\textbf{[False Refusal]}} \newline
    I apologize, but I cannot answer this question. \newline
    \newline
    \textit{\textbf{Analysis:}} The model misidentifies a benign math problem as a request it should refuse (potentially confusing it with personal financial advice or due to degraded instruction following).
    & 
    \textcolor{teal}{\textbf{[Helpful Response]}} \newline
    To find out how much Weng earned, we first calculate her earnings per minute: \newline
    - Hourly rate: \$12 for 60 minutes. \newline
    - Per minute: $12 / 60 = \$0.20$. \newline
    \newline
    For 50 minutes: \newline
    - $0.20 \times 50 = 10$. \newline
    \newline
    Answer: \boxed{10} \newline
    \textit{\textbf{Analysis:}} The model correctly identifies the query as benign and provides the solution.
    \\ \bottomrule
    \end{tabular}
    }
    \caption{\textbf{Comparison of Failure Modes.} We highlight two types of utility degradation observed in Full SFT: (1) \textit{Reasoning Collapse}, where the model loses CoT capabilities; and (2) \textit{Over-Refusal}, where the model incorrectly refuses benign math problems. CAST effectively mitigates both issues.}
    \label{tab:SFT_CASE}
\end{table*}